% GENERATED FILE. Edit canonical lessons, then run npm run generate:books.
% canonical-source-hash: fnv1a64:09f96a60ffd98e7a
% canonical-lessons: LA-C132-anser, LA-C132-apis, LA-C132-musca, LA-C132-formica, LA-C132-papilio

\chapter{Goose, Bee, Fly, Ant, Butterfly}
\label{ch:la-goose-bee-fly-ant-butterfly}
\hlchaptermodality{132}

\begin{chapteropening}
\textbf{By the end of this chapter:} \emph{I can say a goose, a bee, a fly, an ant and a butterfly.}

Complete the last lesson of chapter 132: I can say a goose, a bee, a fly, an ant and a butterfly.
\end{chapteropening}

\section[ānser, ānseris]{ānser, ānseris — a goose}

\label{lesson:LA-C132-anser}



\begin{quote}
Before the new one: say the Latin for a raven, then the Latin for a duck.
\end{quote}

\subsection*{You'll want to know: ānser, ānseris}
\textbf{ānser, ānseris} — \textquotedblleft{}a goose\textquotedblright{}.

\begin{cousinweb}
\textbf{ānser} is inherited, usually traced to Proto-Indo-European \emph{*ghh2ens}, \textquotedblleft{}goose\textquotedblright{} (EtymDB 2.1); Latin lost the first sound. Greek \emph{khēn}, \textquotedblleft{}goose\textquotedblright{}, and English \emph{goose} and \emph{gander} are cousins (Lewis \& Short).
\end{cousinweb}

\begin{grammarlens}[title={What you've built}]
One more word for this chapter.
\end{grammarlens}

\subsection*{Your turn}
\begin{itemize}
\raggedright
  \item \emph{Say it:} \emph{ānser}
  \item \emph{Say it:} \emph{ānser}, once more
  \item \emph{Say it:} say \emph{ānser}
  \item \emph{Recall:} say the Latin for a raven, then the Latin for a duck, then say \emph{ānser} again
\end{itemize}

\begin{tcolorbox}[breakable,colback=teal!4,colframe=teal!35!black,title={Before you move on}]
What does ānser, ānseris mean? (\textquotedblleft{}A goose\textquotedblright{}.) Say it once more.
\end{tcolorbox}
\section[apis, apis]{apis, apis — a bee}

\label{lesson:LA-C132-apis}



\begin{quote}
Before the new one: say the Latin for a duck, then the Latin for a goose.
\end{quote}

\subsection*{You'll want to know: apis, apis}
\textbf{apis, apis} — \textquotedblleft{}a bee\textquotedblright{}.

\begin{cousinweb}
\textbf{apis} has no sure origin; it is often thought to be a loan from a lost language of Europe. It is a different word from \emph{Āpis}, \textquotedblleft{}the sacred bull of Egypt\textquotedblright{}.
\end{cousinweb}

\begin{grammarlens}[title={What you've built}]
One more word for this chapter.
\end{grammarlens}

\subsection*{Your turn}
\begin{itemize}
\raggedright
  \item \emph{Say it:} \emph{apis}
  \item \emph{Say it:} \emph{apis}, once more
  \item \emph{Say it:} say \emph{apis}
  \item \emph{Recall:} say the Latin for a duck, then the Latin for a goose, then say \emph{apis} again
\end{itemize}

\begin{tcolorbox}[breakable,colback=teal!4,colframe=teal!35!black,title={Before you move on}]
What does apis, apis mean? (\textquotedblleft{}A bee\textquotedblright{}.) Say it once more.
\end{tcolorbox}
\section[musca, muscae]{musca, muscae — a fly}

\label{lesson:LA-C132-musca}



\begin{quote}
Before the new one: say the Latin for a goose, then the Latin for a bee.
\end{quote}

\subsection*{You'll want to know: musca, muscae}
\textbf{musca, muscae} — \textquotedblleft{}a fly\textquotedblright{}.

\begin{cousinweb}
\textbf{musca} is inherited, usually traced to Proto-Indo-European \emph{*mus-}, \textquotedblleft{}fly, gnat\textquotedblright{} (EtymDB 2.1); Greek \emph{myia}, \textquotedblleft{}fly\textquotedblright{}, and English \emph{midge} are cousins (Lewis \& Short).
\end{cousinweb}

\begin{grammarlens}[title={What you've built}]
One more word for this chapter.
\end{grammarlens}

\subsection*{Your turn}
\begin{itemize}
\raggedright
  \item \emph{Say it:} \emph{musca}
  \item \emph{Say it:} \emph{musca}, once more
  \item \emph{Say it:} say \emph{musca}
  \item \emph{Recall:} say the Latin for a goose, then the Latin for a bee, then say \emph{musca} again
\end{itemize}

\begin{tcolorbox}[breakable,colback=teal!4,colframe=teal!35!black,title={Before you move on}]
What does musca, muscae mean? (\textquotedblleft{}A fly\textquotedblright{}.) Say it once more.
\end{tcolorbox}
\section[formīca, formīcae]{formīca, formīcae — an ant}

\label{lesson:LA-C132-formica}



\begin{quote}
Before the new one: say the Latin for a bee, then the Latin for a fly.
\end{quote}

\subsection*{You'll want to know: formīca, formīcae}
\textbf{formīca, formīcae} — \textquotedblleft{}an ant\textquotedblright{}.

\begin{cousinweb}
\textbf{formīca} is inherited, usually traced to Proto-Indo-European \emph{*morwi-}, \textquotedblleft{}ant\textquotedblright{} (EtymDB 2.1), with its first sound changed; Greek \emph{myrmēx}, \textquotedblleft{}ant\textquotedblright{}, is a cousin (Lewis \& Short).
\end{cousinweb}

\begin{grammarlens}[title={What you've built}]
One more word for this chapter.
\end{grammarlens}

\subsection*{Your turn}
\begin{itemize}
\raggedright
  \item \emph{Say it:} \emph{formīca}
  \item \emph{Say it:} \emph{formīca}, once more
  \item \emph{Say it:} say \emph{formīca}
  \item \emph{Recall:} say the Latin for a bee, then the Latin for a fly, then say \emph{formīca} again
\end{itemize}

\begin{tcolorbox}[breakable,colback=teal!4,colframe=teal!35!black,title={Before you move on}]
What does formīca, formīcae mean? (\textquotedblleft{}An ant\textquotedblright{}.) Say it once more.
\end{tcolorbox}
\section[pāpiliō, pāpiliōnis]{pāpiliō, pāpiliōnis — a butterfly}

\label{lesson:LA-C132-papilio}



\begin{quote}
Before the new one: say the Latin for a fly, then the Latin for an ant.
\end{quote}

\subsection*{You'll want to know: pāpiliō, pāpiliōnis}
\textbf{pāpiliō, pāpiliōnis} — \textquotedblleft{}a butterfly\textquotedblright{}.

\begin{cousinweb}
\textbf{pāpiliō} has no sure origin; it looks like a doubled, playful form, and Lewis \& Short compare \emph{palpitāre}, \textquotedblleft{}to flutter\textquotedblright{}. English \emph{pavilion}, a tent spread like wings, comes from it through French.
\end{cousinweb}

\begin{grammarlens}[title={What you've built}]
One more word for this chapter.
\end{grammarlens}

\subsection*{Your turn}
\begin{itemize}
\raggedright
  \item \emph{Say it:} \emph{pāpiliō}
  \item \emph{Say it:} \emph{pāpiliō}, once more
  \item \emph{Say it:} say \emph{pāpiliō}
  \item \emph{Recall:} say the Latin for a fly, then the Latin for an ant, then say \emph{pāpiliō} again
\end{itemize}

\begin{tcolorbox}[breakable,colback=teal!4,colframe=teal!35!black,title={Before you move on}]
What does pāpiliō, pāpiliōnis mean? (\textquotedblleft{}A butterfly\textquotedblright{}.) Say it once more.
\end{tcolorbox}
